% =============================================================================
% Chapter 15 -- Recovery Results
% =============================================================================

\chapter{Recovery Results}
\label{ch:recovery}

\section{Overview}
\label{sec:rec:overview}

Phase~2 asks: given a maximally degraded dataset (20\% minority remaining, $\IR \approx 4$), can oversampling restore the performance observed on the fully balanced starting point?
We evaluate seven oversampling methods incrementally as they add minority samples back toward full balance.

The right panels of Figures~\ref{fig:deg_auc}--\ref{fig:deg_spec} in Chapter~\ref{ch:degradation} show the recovery curves for each metric, averaged across all ten datasets and all five classifiers.
Table~\ref{tab:rec_summary} gives mean AUC and F1 at full recovery (100\% of minority restored) per oversampler, compared against the balanced baseline.


% ============================================================================
\section{Recovery Summary Table}
\label{sec:rec:table}

\begin{table}[ht]
\centering
\caption{Mean metric values at full recovery (100\% minority restored via oversampling) vs.\ the balanced baseline.
$\Delta$ = recovery ceiling minus balanced baseline (negative values indicate a residual gap).
Best oversampler per metric in \textbf{bold}.}
\label{tab:rec_summary}
\setlength{\tabcolsep}{5pt}
\begin{tabular}{l ccc ccc}
\toprule
& \multicolumn{3}{c}{AUC} & \multicolumn{3}{c}{F1} \\
\cmidrule(lr){2-4}\cmidrule(lr){5-7}
Oversampler & Balanced & Recovered & $\Delta$ & Balanced & Recovered & $\Delta$ \\
\midrule
No oversampling   & 0.885 & 0.684 & $-0.201$ & 0.850 & 0.541 & $-0.309$ \\
ROS~\cite{fernandez2018learning}        & 0.885 & 0.841 & $-0.044$ & 0.850 & 0.783 & $-0.067$ \\
SMOTE~\cite{chawla2002smote}            & 0.885 & 0.857 & $-0.028$ & 0.850 & 0.802 & $-0.048$ \\
B-SMOTE~\cite{han2005borderline}        & 0.885 & 0.853 & $-0.032$ & 0.850 & 0.811 & $-0.039$ \\
ADASYN~\cite{he2008adasyn}              & 0.885 & 0.860 & $-0.025$ & 0.850 & 0.806 & $-0.044$ \\
KM-SMOTE~\cite{douzas2018improving}     & 0.885 & \textbf{0.871} & $-0.014$ & 0.850 & \textbf{0.827} & $-0.023$ \\
\textbf{GVM-CO}~\cite{hajiannejad2025circular} & 0.885 & \textbf{0.871} & $-0.014$ & 0.850 & 0.824 & $-0.026$ \\
\bottomrule
\end{tabular}
\end{table}


% ============================================================================
\section{RQ3: Full Recovery}
\label{sec:rec:rq3}

No oversampler achieves complete recovery to the balanced baseline for all metrics and all datasets (Table~\ref{tab:rec_summary}).
The gap is most pronounced for Sensitivity: all oversamplers recover to within approximately 90--97\% of the balanced Sensitivity ceiling, but a residual gap persists even when the nominal ratio $\IR = 1$ is re-established via oversampling.

This finding is consistent with the theoretical expectation that synthetic samples carry less information than real minority instances~\cite{chawla2002smote}.
The information content of synthesised examples is bounded by the diversity of the real minority class: when the minority class is small and sparse, the synthesised distribution is an approximation of the original~\cite{blagus2013smote}.
Blagus and Lusa~\cite{blagus2013smote} showed theoretically and empirically that SMOTE cannot recover the full original distribution; our results extend this finding to all seven evaluated oversamplers.

AUC-ROC and Balanced Accuracy show the best recovery (residual gaps of 1.4--4.4\%), consistent with their robustness to threshold choice~\cite{hand2001simple, fernandez2018learning}.
F1-Score and G-Mean show intermediate recovery.
Specificity is largely unaffected by degradation or oversampling (Figure~\ref{fig:deg_spec}), as it measures majority-class accuracy which remains consistent throughout~\cite{he2009learning}.

\paragraph{Residual information loss.}
The residual recovery gap reflects an inherent limitation: synthetic samples carry less information than real minority instances.
When the surviving minority is small, the synthesised distribution is a degraded approximation of the original, regardless of which oversampling method is used.
The theoretical floor on recovery is set by the information content of the 20\% surviving real minority instances; no synthesis method can overcome this information-theoretic constraint.


% ============================================================================
\section{RQ4: Comparative Oversampler Performance}
\label{sec:rec:rq4}

\paragraph{Recovery speed.}
GVM-CO and KM-SMOTE reach 95\% of the balanced AUC within fewer recovery steps than SMOTE or B-SMOTE.
ADASYN shows the fastest initial recovery (first 20 steps) due to adaptive density weighting~\cite{he2008adasyn} but plateaus earlier, as its boundary focus generates samples in majority-adjacent regions.

\paragraph{Recovery ceiling.}
KM-SMOTE~\cite{douzas2018improving} and GVM-CO both achieve mean AUC within 1.4\% of the balanced baseline (Table~\ref{tab:rec_summary}).
Both cluster the minority class before synthesis, preventing interpolation across sub-clusters.
GVM-CO's gravity-biased Von Mises sampling preserves density structure, confirming the mechanistic link between seed quality (Part~III) and recovery (Part~V).

\paragraph{SMOTE variants.}
Borderline-SMOTE recovers Sensitivity more effectively than standard SMOTE but at a slight Specificity cost~\cite{han2005borderline}.
ADASYN shows a similar pattern.

\paragraph{Random Oversampling and lower bound.}
ROS falls further behind at higher ratios due to its lack of diversity.
The no-oversampling baseline's flat recovery curve serves as a lower bound, confirming that all gains are attributable to synthesis rather than evaluation artefacts~\cite{kaufman2012leakage}.

\paragraph{Ranking summary.}
Across AUC and F1, the oversampler ranking is:
GVM-CO $\approx$ KM-SMOTE $>$ ADASYN $>$ SMOTE $>$ B-SMOTE $>$ ROS $\gg$ No-op.
This ranking is consistent with the Part~IV benchmark results (Chapter~\ref{ch:results}), where GVM-CO and LRE-CO outperformed standard SMOTE variants on imbalanced benchmarks.
The convergent evidence from two independent evaluation frameworks (benchmark comparisons in Part~IV and the controlled causal study in Part~V) strengthens confidence in GVM-CO's distributional-fidelity advantage in the regimes studied.
